%!TEX root = ../template.tex
%%%%%%%%%%%%%%%%%%%%%%%%%%%%%%%%%%%%%%%%%%%%%%%%%%%%%%%%%%%%%%%%%%%%
%% abstract-pt.tex
%% NOVA thesis document file
%%
%% Abstract in Portuguese
%%%%%%%%%%%%%%%%%%%%%%%%%%%%%%%%%%%%%%%%%%%%%%%%%%%%%%%%%%%%%%%%%%%%
\typeout{NT FILE abstract-en.tex}%

\glsresetall

A crescente popularidade dos drones tem sido acompanhada por um aumento do número de incursões não intencionais em espaço aéreo restrito, tendo sido documentadas perturbações em aeroportos e outras infraestruturas críticas. Esta tese propõe uma resposta não destrutiva: um enxame de drones que circunda um alvo intruso não cooperativo e forma uma barreira visível, sinalizando ao seu operador que o espaço aéreo é restrito e incentivando o alvo a abandonar a área. Aqui, não cooperativo significa que o alvo não partilha qualquer objetivo de controlo com os perseguidores; não é modelado como um adversário.

A maioria dos métodos de circundamento existentes é planar ou define o circundamento como uma contenção convexa. Este trabalho aborda o caso totalmente tridimensional, no qual os agentes devem ser distribuídos uniformemente em torno do alvo, de modo a que nenhuma direção fique sem cobertura. São desenvolvidas três arquiteturas, com um grau crescente de generalidade. Em primeiro lugar, uma abordagem centralizada posiciona os agentes sobre a esfera unitária utilizando uma malha de Fibonacci e controla os seus movimentos através de um controlador geométrico baseado em Lyapunov, com um controlador de Mellinger para o seguimento de baixo nível. Em segundo lugar, uma abordagem descentralizada permite que os agentes estimem a posição do alvo através de consenso numa rede espacial e se redistribuam utilizando o algoritmo de Lloyd sobre células de Voronoi, com implementações baseadas no invólucro convexo e em semiplanos. Em terceiro lugar, a formação é generalizada para malhas de superfície arbitrárias, incluindo uma variante ponderada pela velocidade e uma configuração de barreira, sendo a geração de trajetórias substituída por um controlo preditivo baseado num modelo (MPC) com polinómios de Bernstein, incorporando restrições de velocidade, aceleração e colisão.

As simulações, realizadas com até 31-agentes, alvos estáticos e móveis e a perda, durante o voo, de vários agentes, demonstram que as formações convergem e se reconstituem em torno do alvo, mantendo-se as distâncias entre agentes acima de um limiar de segurança durante quase toda a simulação. São também apresentadas comparações entre métodos de distribuição de pontos, planeamento de trajetórias e prevenção de colisões. Tanto quanto é do conhecimento do autor, esta combinação específica de geração de formações baseada em superfícies, redistribuição geométrica descentralizada e prevenção de colisões ao nível da trajetória para um alvo móvel não foi estudada de forma sistemática na literatura analisada.

% Palavras-chave do resumo
\keywords{
Circundamento multi-UAV \and
Controlo de formações \and
Controlador baseado em Lyapunov \and
Veículo aéreo não tripulado \and
Veículo terrestre não tripulado \and
Controlo preditivo baseado num modelo \and
Prevenção de colisões \and
Diagramas de Voronoi e algoritmo de Lloyd \and
Planeamento de trajetórias
}

